\section{Higher derived functors}
\label{section-higher-derived}

\noindent
The following simple lemma shows that right derived functors
``move to the right''.

\begin{lemma}
\label{lemma-negative-vanishing}
Let $F : \mathcal{A} \to \mathcal{B}$ be an additive functor
between abelian categories. Let $K^\bullet$ be a complex of $\mathcal{A}$
and $a \in \mathbf{Z}$.
\begin{enumerate}
\item If $H^i(K^\bullet) = 0$ for all $i < a$ and $RF$ is defined at
$K^\bullet$, then $H^i(RF(K^\bullet)) = 0$ for all $i < a$.
\item If $RF$ is defined at $K^\bullet$ and $\tau_{\leq a}K^\bullet$,
then $H^i(RF(\tau_{\leq a}K^\bullet)) = H^i(RF(K^\bullet))$
for all $i \leq a$.
\end{enumerate}
\end{lemma}

\begin{proof}
Assume $K^\bullet$ satisfies the assumptions of (1).
Let $s : K^\bullet \to L^\bullet$ be any quasi-isomorphism.
Then it is also true that $K^\bullet \to \tau_{\geq a}L^\bullet$
is a quasi-isomorphism by our assumption on $K^\bullet$.
Hence in the category $K^\bullet/\text{Qis}^{+}(\mathcal{A})$ the
quasi-isomorphisms $s : K^\bullet \to L^\bullet$ with $L^n = 0$
for $n < a$ are cofinal. From Categories, Lemma
\ref{categories-lemma-cofinal-essentially-constant}
we deduce that $RF$ is the value of the essentially constant
ind-object $F(L^\bullet)$ for these $s$. This means that
$\text{id} : RF(K^\bullet) \to RF(K^\bullet)$ factors through
$F(L^\bullet)$ for some complex $L^\bullet$ with $L^n = 0$
for $n < a$. It follows that $H^i(RF(K^\bullet)) = 0$ for $i < a$.

\medskip\noindent
To prove (2) we use the distinguished triangle
$$
\tau_{\leq a}K^\bullet \to K^\bullet \to \tau_{\geq a + 1}K^\bullet
\to (\tau_{\leq a}K^\bullet)[1]
$$
of Remark \ref{remark-truncation-distinguished-triangle} to conclude
via Lemma \ref{lemma-2-out-of-3-defined} that
$RF$ is defined at $\tau_{\geq a + 1}K^\bullet$ as well and that we have
a distinguished triangle
$$
RF(\tau_{\leq a}K^\bullet) \to RF(K^\bullet) \to RF(\tau_{\geq a + 1}K^\bullet)
\to RF(\tau_{\leq a}K^\bullet)[1]
$$
in $D(\mathcal{B})$. By part (1) we see that $RF(\tau_{\geq a + 1}K^\bullet)$
has vanishing cohomology in degrees $< a + 1$. The long exact cohomology
sequence of this distinguished triangle then shows what we want.
\end{proof}

\begin{definition}
\label{definition-higher-derived-functors}
Let $F : \mathcal{A} \to \mathcal{B}$ be an additive functor
between abelian categories. Assume
$RF : D^{+}(\mathcal{A}) \to D^{+}(\mathcal{B})$ is everywhere
defined. Let $i \in \mathbf{Z}$.
The {\it $i$th right derived functor $R^iF$ of $F$} is the functor
$$
R^iF = H^i \circ RF :
\mathcal{A}
\longrightarrow
\mathcal{B}
$$
\end{definition}

\noindent
The following lemma shows that it really does not make a lot
of sense to take the right derived functor unless the functor
is left exact.

\begin{lemma}
\label{lemma-left-exact-higher-derived}
Let $F : \mathcal{A} \to \mathcal{B}$ be an additive functor
between abelian categories and assume
$RF : D^{+}(\mathcal{A}) \to D^{+}(\mathcal{B})$ is everywhere
defined.
\begin{enumerate}
\item We have $R^iF = 0$ for $i < 0$,
\item $R^0F$ is left exact,
\item the map $F \to R^0F$ is an isomorphism if and
only if $F$ is left exact.
\end{enumerate}
\end{lemma}

\begin{proof}
Let $A$ be an object of $\mathcal{A}$. By Lemma \ref{lemma-negative-vanishing}
we have $H^i(RF(A[0]) = 0$ for $i < 0$. This proves (1).

\medskip\noindent
Let $0 \to A \to B \to C \to 0$ be a short exact sequence of $\mathcal{A}$.
By
Lemma \ref{lemma-derived-canonical-delta-functor}
we obtain a distinguished triangle
$(A[0], B[0], C[0], a, b, c)$ in $D^{+}(\mathcal{A})$.
From the long exact cohomology sequence (and the vanishing for $i < 0$
proved above) we deduce that $0 \to R^0F(A) \to R^0F(B) \to R^0F(C)$
is exact. Hence $R^0F$ is left exact. Of course this also proves that if
$F \to R^0F$ is an isomorphism, then $F$ is left exact.

\medskip\noindent
Assume $F$ is left exact. Recall that $RF(A[0])$ is the value of the
essentially constant system $F(K^\bullet)$ for $s : A[0] \to K^\bullet$
quasi-isomorphisms. It follows that $R^0F(A)$ is the value of the essentially
constant system $H^0(F(K^\bullet))$ for $s : A[0] \to K^\bullet$
quasi-isomorphisms, see Categories, Lemma
\ref{categories-lemma-image-essentially-constant}.
But if $s : A[0] \to K^\bullet$ is a quasi-isomorphism, then
$A[0] \to \tau_{\geq 0}K^\bullet$ is a quasi-isomorphism.
Hence in the category $A[0]/\text{Qis}^{+}(\mathcal{A})$ the
quasi-isomorphisms $s : A[0] \to K^\bullet$ with $K^n = 0$ for $n < 0$
are cofinal. It follows from
Categories, Lemma \ref{categories-lemma-cofinal-essentially-constant}
that we may restrict to such $s$.
Moreover, for such an $s$ the sequence
$$
0 \to A \to K^0 \to K^1
$$
is exact. Since $F$ is left exact we see that
$0 \to F(A) \to F(K^0) \to F(K^1)$ is exact as well.
It follows that $F(A) \to H^0(F(K^\bullet))$ is an isomorphism
and the system is actually constant with value $F(A)$.
We conclude $R^0F = F$ as desired.
\end{proof}

\begin{lemma}
\label{lemma-F-acyclic}
Let $F : \mathcal{A} \to \mathcal{B}$ be an additive functor
between abelian categories and assume
$RF : D^{+}(\mathcal{A}) \to D^{+}(\mathcal{B})$ is everywhere
defined. Let $A$ be an object of $\mathcal{A}$.
\begin{enumerate}
\item $A$ is right acyclic for $F$ if and only if
$F(A) \to R^0F(A)$ is an isomorphism and $R^iF(A) = 0$ for all $i > 0$,
\item if $F$ is left exact, then $A$ is right acyclic for $F$
if and only if $R^iF(A) = 0$ for all $i > 0$.
\end{enumerate}
\end{lemma}

\begin{proof}
If $A$ is right acyclic for $F$, then $RF(A[0]) = F(A)[0]$ and in
particular $F(A) \to R^0F(A)$ is an isomorphism and
$R^iF(A) = 0$ for $i \not = 0$. Conversely, if $F(A) \to R^0F(A)$
is an isomorphism and $R^iF(A) = 0$ for all $i > 0$ then
$F(A[0]) \to RF(A[0])$ is a quasi-isomorphism by
Lemma \ref{lemma-left-exact-higher-derived} part (1)
and hence $A$ is acyclic. If $F$ is left exact then $F = R^0F$, see
Lemma \ref{lemma-left-exact-higher-derived}.
\end{proof}

\begin{lemma}
\label{lemma-F-acyclic-ses}
Let $F : \mathcal{A} \to \mathcal{B}$ be a left exact functor
between abelian categories and assume
$RF : D^{+}(\mathcal{A}) \to D^{+}(\mathcal{B})$ is everywhere
defined. Let $0 \to A \to B \to C \to 0$ be a short exact sequence
of $\mathcal{A}$.
\begin{enumerate}
\item If $A$ and $C$ are right acyclic for $F$ then so is $B$.
\item If $A$ and $B$ are right acyclic for $F$ then so is $C$.
\item If $B$ and $C$ are right acyclic for $F$ and $F(B) \to F(C)$ is
surjective then $A$ is right acyclic for $F$.
\end{enumerate}
In each of the three cases
$$
0 \to F(A) \to F(B) \to F(C) \to 0
$$
is a short exact sequence of $\mathcal{B}$.
\end{lemma}

\begin{proof}
By
Lemma \ref{lemma-derived-canonical-delta-functor}
we obtain a distinguished triangle
$(A[0], B[0], C[0], a, b, c)$ in $D^{+}(\mathcal{A})$.
As $RF$ is an exact functor and since
$R^iF = 0$ for $i < 0$ and $R^0F = F$
(Lemma \ref{lemma-left-exact-higher-derived})
we obtain an exact cohomology sequence
$$
0 \to F(A) \to F(B) \to F(C) \to R^1F(A) \to \ldots
$$
in the abelian category $\mathcal{B}$. Thus the lemma follows from
the characterization of acyclic objects in
Lemma \ref{lemma-F-acyclic}.
\end{proof}

\begin{lemma}
\label{lemma-right-derived-delta-functor}
Let $F : \mathcal{A} \to \mathcal{B}$ be an additive functor
between abelian categories and assume
$RF : D^{+}(\mathcal{A}) \to D^{+}(\mathcal{B})$ is everywhere defined.
\begin{enumerate}
\item The functors $R^iF$, $i \geq 0$ come equipped with a canonical
structure of a $\delta$-functor from $\mathcal{A} \to \mathcal{B}$, see
Homology, Definition \ref{homology-definition-cohomological-delta-functor}.
\item If every object of $\mathcal{A}$ is a subobject of a right
acyclic object for $F$, then $\{R^iF, \delta\}_{i \geq 0}$ is a
universal $\delta$-functor, see
Homology, Definition \ref{homology-definition-universal-delta-functor}.
\end{enumerate}
\end{lemma}

\begin{proof}
The functor $\mathcal{A} \to \text{Comp}^{+}(\mathcal{A})$,
$A \mapsto A[0]$ is exact. The functor
$\text{Comp}^{+}(\mathcal{A}) \to D^{+}(\mathcal{A})$
is a $\delta$-functor, see
Lemma \ref{lemma-derived-canonical-delta-functor}.
The functor $RF : D^{+}(\mathcal{A}) \to D^{+}(\mathcal{B})$ is exact.
Finally, the functor $H^0 : D^{+}(\mathcal{B}) \to \mathcal{B}$
is a homological functor, see
Definition \ref{definition-unbounded-derived-category}.
Hence we get the structure of a $\delta$-functor from
Lemma \ref{lemma-compose-delta-functor-homological}
and
Lemma \ref{lemma-exact-compose-delta-functor}.
Part (2) follows from
Homology, Lemma \ref{homology-lemma-efface-implies-universal}
and the description of acyclics in
Lemma \ref{lemma-F-acyclic}.
\end{proof}

\begin{lemma}[Leray's acyclicity lemma]
\label{lemma-leray-acyclicity}
Let $F : \mathcal{A} \to \mathcal{B}$ be an additive functor
between abelian categories. Let $A^\bullet$ be a bounded below complex
of right $F$-acyclic objects such that $RF$ is defined at
$A^\bullet$\footnote{For example this holds if
$RF : D^{+}(\mathcal{A}) \to D^{+}(\mathcal{B})$ is everywhere defined.}.
The canonical map
$$
F(A^\bullet) \longrightarrow RF(A^\bullet)
$$
is an isomorphism in $D^{+}(\mathcal{B})$, i.e., $A^\bullet$ computes
$RF$.
\end{lemma}

\begin{proof}
Let $A^\bullet$ be a bounded complex of right $F$-acyclic objects.
We claim that $RF$ is defined at $A^\bullet$ and that
$F(A^\bullet) \to RF(A^\bullet)$ is an isomorphism in $D^+(\mathcal{B})$.
Namely, it holds for complexes with at most one nonzero right $F$-acyclic
object by Definition \ref{definition-derived-functor}.
Next, suppose that
$A^n = 0$ for $n \not \in [a, b]$. Using the ``stupid'' truncations we obtain
a termwise split short exact sequence of complexes
$$
0 \to \sigma_{\geq a + 1} A^\bullet \to A^\bullet \to
\sigma_{\leq a} A^\bullet \to 0
$$
see
Homology, Section \ref{homology-section-truncations}.
Thus a distinguished triangle
$(\sigma_{\geq a + 1} A^\bullet, A^\bullet, \sigma_{\leq a} A^\bullet)$.
By induction hypothesis $RF$ is defined for the two outer complexes
and these complexes compute $RF$. Then the same is true for the middle one
by Lemma \ref{lemma-2-out-of-3-computes}.

\medskip\noindent
Suppose that $A^\bullet$ is a bounded below complex of acyclic objects
such that $RF$ is defined at $A^\bullet$.
To show that $F(A^\bullet) \to RF(A^\bullet)$
is an isomorphism in $D^{+}(\mathcal{B})$
it suffices to show that $H^i(F(A^\bullet)) \to H^i(RF(A^\bullet))$
is an isomorphism for
all $i$. Pick $i$. Consider the termwise split short exact sequence of
complexes
$$
0 \to \sigma_{\geq i + 2} A^\bullet \to A^\bullet \to
\sigma_{\leq i + 1} A^\bullet \to 0.
$$
Note that this induces a termwise split short exact sequence
$$
0 \to \sigma_{\geq i + 2} F(A^\bullet) \to F(A^\bullet) \to
\sigma_{\leq i + 1} F(A^\bullet) \to 0.
$$
Hence we get distinguished triangles
$$
(\sigma_{\geq i + 2} A^\bullet, A^\bullet,
\sigma_{\leq i + 1} A^\bullet)
\quad\text{and}\quad
(\sigma_{\geq i + 2} F(A^\bullet), F(A^\bullet),
\sigma_{\leq i + 1} F(A^\bullet))
$$
Since $RF$ is defined at $A^\bullet$ (by assumption)
and at $\sigma_{\leq i + 1}A^\bullet$ (by the first paragraph)
we see that $RF$ is defined at $\sigma_{\geq i + 2}A^\bullet$
and we get a distinguished triangle
$$
(RF(\sigma_{\geq i + 2} A^\bullet), RF(A^\bullet),
RF(\sigma_{\leq i + 1} A^\bullet))
$$
See Lemma \ref{lemma-2-out-of-3-defined}.
Using these distinguished triangles we obtain a map of exact sequences
$$
\xymatrix{
H^i(\sigma_{\geq i + 2} F(A^\bullet)) \ar[r] \ar[d] &
H^i(F(A^\bullet)) \ar[r] \ar[d]^\alpha &
H^i(\sigma_{\leq i + 1} F(A^\bullet)) \ar[r] \ar[d]^\beta &
H^{i + 1}(\sigma_{\geq i + 2} F(A^\bullet)) \ar[d] \\
H^i(RF(\sigma_{\geq i + 2} A^\bullet)) \ar[r] &
H^i(RF(A^\bullet)) \ar[r] &
H^i(RF(\sigma_{\leq i + 1} A^\bullet)) \ar[r] &
H^{i + 1}(RF(\sigma_{\geq i + 2} A^\bullet))
}
$$
By the results of the first paragraph the map $\beta$ is an isomorphism.
By inspection the objects on the upper left and the upper right
are zero. Hence to finish the proof it suffices to show that
$H^i(RF(\sigma_{\geq i + 2} A^\bullet)) = 0$ and
$H^{i + 1}(RF(\sigma_{\geq i + 2} A^\bullet)) = 0$.
This follows immediately from
Lemma \ref{lemma-negative-vanishing}.
\end{proof}

\begin{proposition}
\label{proposition-enough-acyclics}
\begin{slogan}
A functor on Abelian categories is extended to the (bounded below or above)
derived category by resolving with a complex that is acyclic for that functor.
\end{slogan}
Let $F : \mathcal{A} \to \mathcal{B}$ be an additive functor of
abelian categories.
\begin{enumerate}
\item If every object of $\mathcal{A}$ injects into an object acyclic
for $RF$, then $RF$ is defined on all of $K^{+}(\mathcal{A})$
and we obtain an exact functor
$$
RF : D^{+}(\mathcal{A}) \longrightarrow D^{+}(\mathcal{B})
$$
see (\ref{equation-everywhere}). Moreover, any bounded below complex
$A^\bullet$ whose terms are acyclic for $RF$ computes $RF$.
\item If every object of $\mathcal{A}$ is quotient of
an object acyclic for $LF$, then $LF$ is defined on all of
$K^{-}(\mathcal{A})$ and we obtain an exact functor
$$
LF : D^{-}(\mathcal{A}) \longrightarrow D^{-}(\mathcal{B})
$$
see (\ref{equation-everywhere}). Moreover, any bounded above complex
$A^\bullet$ whose terms are acyclic for $LF$ computes $LF$.
\end{enumerate}
\end{proposition}

\begin{proof}
Assume every object of $\mathcal{A}$ injects into an object acyclic
for $RF$. Let $\mathcal{I}$ be the set of objects acyclic for $RF$.
Let $K^\bullet$ be a bounded below complex in $\mathcal{A}$. By
Lemma \ref{lemma-subcategory-right-resolution}
there exists a quasi-isomorphism $\alpha : K^\bullet \to I^\bullet$ with
$I^\bullet$ bounded below and $I^n \in \mathcal{I}$. Hence in order to
prove (1) it suffices to show that
$F(I^\bullet) \to F((I')^\bullet)$ is a quasi-isomorphism when
$s : I^\bullet \to (I')^\bullet$ is a quasi-isomorphism of bounded
below complexes of objects from $\mathcal{I}$, see
Lemma \ref{lemma-find-existence-computes}.
Note that the cone $C(s)^\bullet$ is an acyclic bounded below complex
all of whose terms are in $\mathcal{I}$.
Hence it suffices to show: given an acyclic bounded below complex
$I^\bullet$ all of whose terms are in $\mathcal{I}$ the complex
$F(I^\bullet)$ is acyclic.

\medskip\noindent
Say $I^n = 0$ for $n < n_0$. Setting $J^n = \Im(d^n)$ we break
$I^\bullet$ into short exact sequences
$0 \to J^n \to I^{n + 1} \to J^{n + 1} \to 0$
for $n \geq n_0$. These sequences induce distinguished triangles
$(J^n, I^{n + 1}, J^{n + 1})$ in $D^+(\mathcal{A})$ by
Lemma \ref{lemma-derived-canonical-delta-functor}.
For each $k \in \mathbf{Z}$ denote $H_k$ the assertion:
For all $n \leq k$ the object $J^n$ is in $\mathcal{I}$.
Then $H_k$ holds trivially for $k < n_0$. If $H_n$ holds,
then Lemma \ref{lemma-2-out-of-3-computes} shows that
$J^{n + 1}$ is in $\mathcal{I}$ and we have $H_{n + 1}$.
By Proposition \ref{proposition-derived-functor} we have a
distinguished triangle $(RF(J^n), RF(I^{n + 1}), RF(J^{n + 1}))$.
Since $J^n, I^{n + 1}, J^{n + 1}$ are in $\mathcal{I}$
the long exact cohomology sequence
(\ref{equation-long-exact-cohomology-sequence-D})
associated to this distinguished triangle collapses
to an exact sequence
$$
0 \to F(J^n) \to F(I^{n + 1}) \to F(J^{n + 1}) \to 0
$$
This in turn proves that $F(I^\bullet)$ is exact.

\medskip\noindent
The proof in the case of $LF$ is dual.
\end{proof}

\begin{lemma}
\label{lemma-right-derived-exact-functor}
Let $F : \mathcal{A} \to \mathcal{B}$ be an exact functor of
abelian categories. Then
\begin{enumerate}
\item every object of $\mathcal{A}$ is right acyclic for $F$,
\item $RF : D^{+}(\mathcal{A}) \to D^{+}(\mathcal{B})$ is everywhere defined,
\item $RF : D(\mathcal{A}) \to D(\mathcal{B})$ is everywhere defined,
\item every complex computes $RF$, in other words, the canonical
map $F(K^\bullet) \to RF(K^\bullet)$ is an isomorphism for all complexes, and
\item $R^iF = 0$ for $i \not = 0$.
\end{enumerate}
\end{lemma}

\begin{proof}
This is true because $F$ transforms acyclic complexes into acyclic complexes
and quasi-isomorphisms into quasi-isomorphisms. Details omitted.
\end{proof}









